\section*{第 \ref{ch:b3:advanced-geochronology} 章——高级地质年代学}

\begin{solution}{exo:b3:advanced-geochronology:1}
$t_{206} = \ln(1.04500)/\lambda_{238} = \qty{283.8}{\Ma}$; $t_{207} =
\ln(1.3300)/\lambda_{235} = \qty{289.6}{\Ma}$，两者相差
\qty{6}{\Ma}：该锆石不谐和。其 $t_{207}$ 较老，因此数据点位于
\omterm{def:b3:advanced-geochronology:concordia}{谐和线}下方，处于铅丢失的一侧。
\end{solution}

\begin{solution}{exo:b3:advanced-geochronology:2}
$J = (e^{5.543\times10^{-10}\times2.8201\times10^7} - 1)/1.3129 =
0.015754/1.3129 = 0.01200$。然后 $t = \ln(1 + 0.01200\times10.000)/\lambda =
\ln 1.12/(5.543\times10^{-10}) = \qty{204.5}{\Ma}$.
\end{solution}

\begin{solution}{exo:b3:advanced-geochronology:3}
$T_{\mathrm{CHUR}} = \ln\bigl(1 + (0.51180 - 0.512638)/(0.1200 -
0.1966)\bigr)/\lambda = \ln(1.010940)/\lambda = \qty{1.67}{\Ga}$;
$T_{\mathrm{DM}} = \ln\bigl(1 + (0.51180 - 0.51315)/(0.1200 - 0.2137)\bigr)/\lambda
= \ln(1.014408)/\lambda = \qty{2.19}{\Ga}$.
\end{solution}

\begin{solution}{exo:b3:advanced-geochronology:4}
不计衰变时，$t = N/P_0 = \num{2.0e5}/4.96 = \qty{40.3}{ka}$。考虑衰变时，
$t = -\ln(1 - N\lambda/P_0)/\lambda = -\ln(1 - 0.02014)/\lambda =
\qty{40.7}{ka}$，年龄偏大 \qty{1}{\percent}。
\end{solution}

\begin{solution}{exo:b3:advanced-geochronology:5}
对较小的 $t$，$x \approx \lambda_{235}t$，$y \approx \lambda_{238}t$，于是
$y \approx (\lambda_{238}/\lambda_{235})x$，斜率 $0.1575$。若铅
近期才丢失，则 $t_2 \approx 0$，\omterm{def:b3:advanced-geochronology:concordia}{下交点}为 $C(0)$，
即原点。沿过原点的直线，$y/x$ 为常数，
$^{207}\mathrm{Pb}^*/^{206}\mathrm{Pb}^* = x/(137.8\,y)$ 亦然，每颗晶体给出
\omterm{def:b3:advanced-geochronology:concordia}{上交点}的铅–铅年龄，即形成年龄。
\end{solution}

\begin{solution}{exo:b3:advanced-geochronology:6}
权重为 $1.5625$、$1$、$1.2346$、$2.0408$，$W = 5.838$；$\bar x =
\qty{296.09}{\Ma}$，$\sigma = 1/\sqrt W = \qty{0.41}{\Ma}$。加权
平方偏差之和为 $2.32$；取 $f = 3$，$\mathrm{MSWD} = 0.77$，小于
$1 + 2\sqrt{2/3} = 2.63$，年龄相符。
\end{solution}

\begin{solution}{exo:b3:advanced-geochronology:7}
$\lambda\zeta g\rho_d\rho_s/\rho_i = 1.55125\times10^{-10}\times350\times0.5
\times\num{1.10e6}\times0.5 = 0.014931$; $t = \ln(1.014931)/\lambda =
\qty{95.5}{\Ma}$。若岩石快速冷却，该年龄记录其穿过
磷灰石部分退火带，约 60 至 \qty{120}{\celsius}；若岩石
缓慢冷却，则为混合年龄，须由径迹长度判别。
\end{solution}

\begin{solution}{exo:b3:advanced-geochronology:8}
每克：铀 $20\times10^{-6}/238.05\times N_A = \num{5.06e16}$ 个原子，其中
铀-238 为 \num{5.02e16}，铀-235 为 \num{3.64e14}；钍
\num{1.038e17}, $\lambda_{232} = \ln 2/\num{1.40e10} = \qty{4.95e-11}{\text{年}^{-1}}$;
氦 $10^{-9}N_A = \num{6.02e14}$。\omterm{def:b3:advanced-geochronology:cosmogenic}{产生率}为
$8\lambda_{238}{}^{238}\mathrm U + 7\lambda_{235}{}^{235}\mathrm U +
6\lambda_{232}\mathrm{Th} = (6.23 + 0.25 + 3.08)\times10^7 =
\num{9.57e7}\text{ 个原子}\,\unit{g^{-1}.\text{年}^{-1}}$，则 $t = \num{6.02e14}/\num{9.57e7} =
\qty{6.29}{Ma}$。校正后：$6.29/0.75 = \qty{8.39}{\Ma}$。
\end{solution}

\begin{solution}{exo:b3:advanced-geochronology:9}
$P_0/\lambda = 4.96/\num{4.994e-7} = \num{9.93e6}\text{ 个原子}\,\unit{g^{-1}}$，达到
\qty{90}{\percent} 需时 $\ln 10/\lambda = \qty{4.6}{Ma}$。取
\qty{10}{m.Ma^{-1}} $= \qty{1e-3}{cm.\text{年}^{-1}}$: $\rho\varepsilon/\Lambda =
2.7\times10^{-3}/160 = \qty{1.69e-5}{\text{年}^{-1}}$, $\mu = \qty{1.74e-5}{\text{年}^{-1}}$
则 $N = 4.96/\num{1.74e-5} = \num{2.85e5}\text{ 个原子}\,\unit{g^{-1}}$，少 35 倍。
\end{solution}

\begin{solution}{exo:b3:advanced-geochronology:10}
记 $W = w_1 + w_2$，$x_1 - \bar x = w_2(x_1 - x_2)/W$，且 $x_2 - \bar x =
-w_1(x_1 - x_2)/W$，于是 $\sum w_i(x_i - \bar x)^2 = (x_1 - x_2)^2(w_1w_2^2 +
w_2w_1^2)/W^2 = (x_1 - x_2)^2w_1w_2/W = (x_1 - x_2)^2/(\sigma_1^2 +
\sigma_2^2)$；取 $f = 1$ 时即为 MSWD。若两者测得相同年龄，
$x_1 - x_2$ 的均值为 0，方差为 $\sigma_1^2 + \sigma_2^2$，因此
其平方的期望与该方差之比为 1。
\end{solution}

\begin{solution}{exo:b3:advanced-geochronology:11}
$C(2700) = (13.283, 0.52019)$, $C(500) = (0.63628, 0.08065)$. $f = (0.52019 -
0.300)/(0.52019 - 0.08065) = 0.501$，即其中一半为铅。$x = 0.499\times13.283 +
0.501\times0.63628 = 6.948$. $t_{206} = \qty{1691}{\Ma}$, $t_{207} =
\qty{2105}{\Ma}$；比值 $x/(137.8\,y) = 0.1681$ 经迭代，得
铅–铅年龄为 \qty{2537}{\Ma}。三个数值均无意义：它们
是两次事件的混合，只有交点才能给出其年代。铅–铅
年龄是结晶年龄的下限。
\end{solution}

\begin{solution}{exo:b3:advanced-geochronology:12}
无剥蚀：$t = -\ln(1 - \num{1.0e6}\times\num{4.994e-7}/4.96)/\lambda =
\qty{212}{ka}$。稳态：$\rho\varepsilon/\Lambda = P_0/N - \lambda =
\qty{4.46e-6}{\text{年}^{-1}}$, $\varepsilon = \qty{2.64e-4}{cm.\text{年}^{-1}} =
\qty{2.6}{m.Ma^{-1}}$。对低于此值的任意 $\varepsilon$，$N =
P_0(1 - e^{-\mu t})/\mu$ 确定一个 $t$，其取值自 \qty{212}{ka} 起
直至无穷大：一次测量含有两个未知数。另取半衰期不同的核素，
如同一石英中的铝-26，或岩芯某深度处的浓度，
即可将二者分开。
\end{solution}

\begin{solution}{pb:b3:advanced-geochronology:1}
\textbf{1.} $t_{206}$ / $t_{207}$，以 \unit{\Ma} 计：Z1 301.9 / 302.2，Z2 284.1 /
286.5，Z3 265.7 / 268.5，Z4 246.1 / 250.1，Z5 225.1 / 229.4。误差
\qty{0.2}{\percent}（作用于 $y$ 上）约合 \qty{0.6}{\Ma}：Z1 在误差范围内谐和，
其余各点不谐和，自 Z2 至 Z5 程度渐增。

\textbf{2.} 锆石在之后某一次事件中丢失的铅
比例逐渐增大：它们沿弦线越来越靠下。弦线位于
凹的\omterm{def:b3:advanced-geochronology:concordia}{谐和线}下方，$x$ 较大，即对给定的 $y$ 而言：$t_{207} > t_{206}$。

\textbf{3.} Z1：$X = 1/0.04794 = 20.859$，$Y = 0.3467/(137.8\times0.04794) =
0.05248$。Z5：$X = 28.145$，$Y = 0.05178$。

\textbf{4.} 权重为 $1.000$、$1.235$、$1.563$、$1.563$、$2.041$（以
$10^8$ 为单位）；$\bar x = 0.29402$，$\bar y = 0.040892$，$S_{xx} = \num{7.845e5}$。

\textbf{5.} $b = 0.13260$, $a = \bar y - b\bar x = 0.001905$.

\textbf{6.} 在 \qty{310}{\Ma} 处，$\mathrm dy/\mathrm dx = 0.1575\,
e^{(\lambda_{238} - \lambda_{235})t} = 0.1218$。\omterm{def:b3:advanced-geochronology:concordia}{不一致线}仅
比\omterm{def:b3:advanced-geochronology:concordia}{谐和线}陡 \qty{9}{\percent}，二者以小角度相交，
故\omterm{def:b3:advanced-geochronology:concordia}{上交点}的精度较低。

\textbf{7.} 用牛顿法从 \qty{310}{\Ma} 出发：$t = \qty{311.2}{\Ma}$。

\textbf{8.} 从 \qty{100}{\Ma} 出发：$t = \qty{99.7}{\Ma}$。

\textbf{9.} $(y_i - a - bx_i)/\sigma_i = 0.64$, $-0.91$, $0.35$, $-0.24$,
$0.17$；$\mathrm{MSWD} = 1.435/3 = 0.48$，低于界限 $2.63$（$f = 3$）：
这些锆石符合同一条直线。

\textbf{10.} $s_b = 1/\sqrt{\num{7.845e5}} = 0.00113$.

\textbf{11.} 取 $b - 2s_b = 0.13034$ 时，\omterm{def:b3:advanced-geochronology:concordia}{上交点}为
\qty{322.9}{\Ma}；取 $b + 2s_b = 0.13485$ 时为 \qty{302.6}{\Ma}。年龄为
$311 \pm \qty{10}{\Ma}$（两倍标准差）。

\textbf{12.} $f = (x(311.2) - x_i)/(x(311.2) - x(99.7))$：Z1 为 \qty{4.7}{\percent}，
Z5 为 \qty{41}{\percent}。

\textbf{13.} $J = 0.015754/1.3129 = 0.0119999$, $0.01200$.

\textbf{14.} 第 1 步：\qty{150.0}{\Ma}；第 4 步：\qty{296.6}{\Ma}；第 10 步：
\qty{295.0}{\Ma}。

\textbf{15.} 第 4 至 10 步：连续七步，共占 \qty{82}{\percent} 的
氩-39，无明显趋势。第 1 至 3 步呈阶梯状上升。

\textbf{16.} 误差相等：平均值为 $2072.9/7 = \qty{296.1}{\Ma}$，$\sigma =
1.1/\sqrt7 = \qty{0.42}{\Ma}$；偏差平方和为 $3.83$，
$\mathrm{MSWD} = 3.83/(1.21\times6) = 0.53$，低于 $2.15$。

\textbf{17.} $\bar R = \sum f_iR_i = 14.270$，$t = \qty{285.2}{\Ma}$：
边缘丢失的氩在前几步释放，拉低了平均值，
正如它会拉低钾–氩年龄一样。

\textbf{18.} $(700 - 375)/(311.2 - 296.1) = \qty{21.5}{\celsius.Ma^{-1}}$
（因锆石年龄的不确定度为 $\pm\qty{10}{\Ma}$，故取值范围较宽）。

\textbf{19.} 约 \qty{100}{\Ma} 时，花岗岩被埋在其中生代盖层之下，
加热或流体循环使锆石丢失了铅。
黑云母边缘丢失氩，阶梯谱记录了这一过程；其
最低一步 \qty{150}{\Ma} 是老氩与重置氩的混合。\omterm{def:b3:advanced-geochronology:arar}{坪年龄}
得以保留：黑云母受热不足，未丢失全部氩。

\textbf{20.} 海拔对年龄的最小二乘斜率为
\qty{20.1}{m.Ma^{-1}}。

\textbf{21.} $(65 - 10)/30 = \qty{1.83}{km}$；自 \qty{43.1}{\Ma} 以来，平均速率为
$1830/43.1 = \qty{42}{m.Ma^{-1}}$。

\textbf{22.} $\rho\varepsilon/\Lambda = P_0/N - \lambda = \num{3.415e-5} -
\num{0.050e-5} = \qty{3.365e-5}{\text{年}^{-1}}$; $\varepsilon = \num{3.365e-5}\times
160/2.7 = \qty{1.99e-3}{cm.\text{年}^{-1}} = \qty{19.9}{m.Ma^{-1}}$.

\textbf{23.} $P_0\Lambda/(\rho N) = \qty{20.2}{m.Ma^{-1}}$，多出 \qty{1.5}{\percent}
：此速率下衰变几乎不起作用。

\textbf{24.} $\Lambda/(\rho\varepsilon) = 160/(2.7\times\num{1.99e-3}) =
\qty{30}{ka}$。花岗岩在 52 至
\qty{43}{\Ma} 之间约以 \qty{20}{m.Ma^{-1}} 的速率抬升，此后平均为 \qty{42}{m.Ma^{-1}}，
现今侵蚀速率为 \qty{20}{m.Ma^{-1}}；
新生代某一时段剥露曾更快。

\textbf{25.} 瓦尔登花岗岩结晶于距今
\textbf{\qty[text-series-to-math=true, separate-uncertainty=true]{311 \pm 10}{\Ma}}（\omterm{def:b3:advanced-geochronology:concordia}{上交点}，
两倍标准差）。
\end{solution}
